\clearpage
\largerpage
\sscp{Innovations among the vowels}{08-01-01}
As noted in \trf{08-01},
for \tsc{Taliabu}, *-ay > \it{e}, *-aw > \it{o},
and *e > \it{o}.
The first two are diagnostic of \tsc{Celebic} languages,
and the third of \tsc{Eastern Celebic} \citep{Mead2003a,Mead2003b}.
At the same time,
these three sound changes systematically set these languages
apart from all other subgroups in geographical central Maluku,
including \tsc{Sula-Buru}, \tsc{Ambon-Seram}, and \tsc{Seram-Tanimbar-Bomberai}.
The correspondences *-ay > \it{e}
and *-aw > \it{o} are also found in some \tsc{Timor-Babar}
and \tsc{Bima-Lembata}
languages.
\trf{08-04} shows reflexes of PMP *-ay > \it{e}
and *-aw > \it{o}
and \trf{08-05} shows reflexes of PMP *ə.

\begin{table}\tcp{\tsc{Taliabu} and nearby languages *-ay and *-aw }{08-04}
%\begin{adjustbox}{width=\textwidth}
	\begin{threeparttable}\st{2.5pt}
	\begin{tabular}{llll|lll}\lsptoprule
local gl.	&	die	&	liver	&	 chin, jaw	&	rafter	&	sister (m.s.)	&	steal\\
PMP	&	*mat\tbr{ay}	&	*qat\tbr{ay}	&	*qaz\tbr{ay}	&	*kas\tbr{aw}	&	*bət\tbr{aw}	&	**nak\tbr{aw}\\\midrule
Saluan	&	{\hr}mat\tbr{e}	&	{\hr}at\tbr{e}	&	{\hr}aʤ\tbr{e}	&	{\hr}kas\tbr{o}ʔ	&	\cg{(utus)}	&	\cg{(mombuliŋ)}\\
Batui	&	{\hr}mat\tbr{e}	&	{\hr}at\tbr{e}	&	{\hr}aʤ\tbr{e}	&	{\hr}kas\tbr{o}ʔ	&	\cg{(utus)}	&	{\hr}mi-nak\tbr{o}ʔ\\
Balantak	&	{\hr}pat\tbr{e}	&	{\hr}at\tbr{e}	&	{\hr}as\tbr{i}	&	{\hr}kas\tbr{o}	&	\cg{(utus)}	&	\cg{(ma-maŋan)}\\
Banggai\su{†}	&	{\hr}mat\tbr{e}	&	{\hr}at\tbr{e}	&	{\hr}ad\tbr{e}	&	{\hr}kaso	&		&	\cg{(mbonikuk)}\\
\dshmidrule{7}
Soboyo	&	{\hr}mat\tbr{e}	&	{\hr}n-at\tbr{e}	&	{\hr}n-ad\tbr{e}\su{Δ}	&	{\hr}kas\tbr{o}	&	{\hr}fot\tbr{o}	&	{\hr}paka-nak\tbr{o}\\
Taliabu	&	{\hr}mat\tbr{e}	&	{\hr}at\tbr{e}	&	{\hr}ad\tbr{e}\su{Δ}	&	{\hr}kas\tbr{o}	&	{\hr}fot\tbr{o}	&	{\hr}paka-nak\tbr{o}\\
Kadai\su{‡}	&	{\hr}ka-mat\tbr{e}	&	{\hr}n-at\tbr{e}	&	{\hr}ad\tbr{e}\su{Δ}	&		&	{\hr}fot\tbr{o}	&	\\\midrule
Mangole	&	{\hr}mat\tbb{a}	&		&		&		&		&	{\hr}bi-nak\tbb{a}\\
Sanana	&	{\hr}mat\tbb{a}	&	\cg{(kila)}	&	{\hr}n-ay\tbb{a}	&	{\hr}kah\tbb{a}	&	{\hr}fet\tbb{a}	&	{\hr}bil-nak\tbb{a}\\
Lisela	&	{\hr}mata-ʔ	&	\cg{(lale-n)}	&	{\hr}a\tbb{a}-n	&	{\hr}kas\tbb{a}	&	{\hr}ɸet\tbb{a}	&	\\
Buru	&	{\hr}mat\tbb{a}	&	{\hr}at\tbb{a}-n\su{ǁ}	&	{\hr}a\tbb{a}-n	&	{\hr}kas\tbb{a}	&	{\hr}fet\tbb{a}	&	{\hr}eʔnak\tbb{a}\\
Ambelau	&	{\hr}em-mar\tbb{a}	&		&	{\hr}al\tbb{a}-ni	&		&	{\hr}ber\tbb{a}	&	\\\midrule
Kayeli 	&	{\hr}en-mat\tbb{a}	&	\cg{(lale-n)}	&	{\hr}a\tbb{a}-ni (<SB)	&		&	{\hr}bet\tbb{a}	&	\\
		\lspbottomrule
	\end{tabular}
		\begin{tablenotes}\tnsize
			\item[†]	{Additional Banggai correspondences supporting *-aw > \it{o}
								and *-ay > \it{e} include: *ma-linaw ‘calm (water)’ > \it{mo-lino},
								*babaw ‘above’ > \it{babo}; *bəRsay ‘paddle’ > \it{bose}.}
			\item[‡]	{Kadai *-aw > \it{o} is also supported by *ka-labaw > \it{kalafo} ‘rat, mouse’.}
			\item[Δ]	{Soboyo \it{n-ade},
								Taliabu, Kadai \it{ade} ‘chin’ are probable loans from Banggai
								given *z > \it{d} in Banggai,
								but *z > \it{y} in \tsc{Taliabu} (see \trf{08-08}).}
			\item[{ǁ}]	{Buru \it{ata-n} is ‘deer spleen’.}
		\end{tablenotes}
	\end{threeparttable}
%\end{adjustbox}
\end{table}

\begin{table}\tcp{\tsc{Taliabu} and nearby languages *ə}{08-05}
\begin{adjustbox}{width=\textwidth}
	\begin{threeparttable}\st{2pt}
	\begin{tabular}{llllllll}\lsptoprule
*gloss	&	 gall, bile	&	new	&	stand	&	thatch	&	sniff	&	charcoal	&	seize\\
PMP	&	*qap\tbr{ə}ju	&	*baq\tbr{ə}Ru	&	*k\tbr{ə}d\tbr{ə}ŋ	&	*qat\tbr{ə}p	&	*haj\tbr{ə}k	&	*qaj\tbr{ə}ŋ	&	*dak\tbr{ə}p\\ \midrule
Saluan	&	{\hr}p\tbr{o}uʔ	&	{\hr}buʔ\tbr{o}u	&	\cg{(tinʤo)}	&	{\hr}at\tbr{o}p	&	{\hr}o\tbr{o}k	&	\cg{(buhiŋ)}	&	\\
Balantak	&	(s)op\tbr{o}yuʔ	&	{\hr}uʔ\tbr{u}ru	&	{\hr}-kerer	&	{\hr}at\tbr{o}p	&	{\hr}o\tbr{o}k	&	{\hr}alay\tbr{o}ŋ	&	{\hr}-rak\tbr{o}p\\
Banggai	&	\cg{(sopot)}	&	\cg{(soodo)}	&		&		&	{\hr}oy\tbr{o}ki\su{†}	&	{\hr}uy\tbr{o}ŋ	&	\\
\noalign{\vskip-1.5ex}
\multicolumn{8}{@{}c@{}}{\dashedline{130mm}} \\
\noalign{\vskip-0.5ex}
Soboyo	&	{\hr}n-\tbr{o}yo-in	&	{\hr}f\tbr{o}hu	&	{\hr}k\tbr{o}h\tbr{o}-in	&	{\hr}at\tbr{o}	&	{\hr}hay\tbr{o}ʔ	&	{\hr}ay\tbr{o}ŋ	&	{\hr}hak\tbr{o}\\
Taliabu	&	{\hr}n-\tbr{o}yu-ɲ	&	{\hr}f\tbr{o}hu	&	{\hr}k\tbr{o}h\tbr{o}ŋ	&	{\hr}at\tbr{o}	&	{\hr}hay\tbr{o}k	&	{\hr}ay\tbr{o}ŋ	&	{\hr}hak\tbr{o}\\
Kadai	&	{\hr}n-\tbr{o}yu	&	{\hr}f\tbr{o}hu	&		&		&	{\hr}hay\tbr{o}	&	{\hr}ay\tbr{o}ŋ	&	{\hr}hak\tbr{o} \it{‘carry’}\\ \midrule
Mangole	&	{\hr}n-p\tbb{e}u	&	{\hr}f\tbb{e}u	&	{\hr}k\tbb{e}li, g\tbb{e}li	&	\cg{(saŋapetu)}	&	\cg{(suun)}	&		&	\\
Sanana	&	{\hr}m-p\tbb{e}u	&	{\hr}f\tbb{e}u	&	{\hr}g\tbb{e}hi	&	\cg{(sanapet)}	&	\cg{(muhi)}	&		&	{\hr}hak kot\\
Lisela	&	{\hr}p\tbb{e}u-n	&	{\hr}ɸ\tbb{e}hu	&	{\hr}k\tbb{e}re	&	\cg{(aba-t)}	&		&		&	\cg{(pogo)}\\
Buru	&	{\hr}p\tbb{e}u-n	&	{\hr}f\tbb{e}hu-t	&	{\hr}k\tbb{e}re-k	&	{\hr}at\tbb{e}-t	&	{\hr}ma\tbb{e}-k	&	\cg{(ime-n)}	&	{\hr}ef-rak\tbb{e}\\
Ambelau	&	{\hr}f\tbb{e}o-ni	&	{\hr}b\tbb{i}hu	&	\cg{(en-βati)}	&	\cg{(rae, lau)}	&		&		&	\\ \midrule
Kayeli 	&		&	{\hr}b\tbb{e}lo	&	\cg{(sabele)}	&	\mc{2}{l}{{\hr}atap (<Mly.)}		&	\\
		\lspbottomrule
	\end{tabular}
		\begin{tablenotes}\tnsize
			\item[†] {Banggai \it{oyoki} means ‘kiss’.}
		\end{tablenotes}
	\end{threeparttable}
\end{adjustbox}
\end{table}

A notable example of *ə > \it{o} is *wahiyəR ‘water’ > Soboyo, Kadai \it{wayo} ‘water’.
The changes of *y = \it{y}
and *ə > \it{o},
show that this cannot derive from a form like *wahiR through
systematic sound changes, so a reconstruction with
*ə like that from \citet{Wolff2010}
or a secondary **wayəR is necessary to account for not only the
\tsc{Taliabu} data, but also most of the \tsc{Sula-Buru} and \tsc{Ambon-Seram} data.

\clearpage
A conditioned exception to PMP *ə > \it{o} is found in many words
that end in *-əj{\#} > \it{e,
o}, as shown in \trf{08-06}.
In \tsc{Sula-Buru}, because final *j{\#} > {\0}
and *ə > \it{e},
\tsc{Sula-Buru} languages also have /e/ in these words,
albeit through different processes.
The presence of *j,
with subsequent *j > **y,
needs to be reconstructed for the \tsc{Taliabu} languages to
account for words that have PMP *əj{\#} > \it{e},
rather than \it{o}, similar to what \citet[76--77]{Mead2003b}
argues for \tsc{Saluan-Banggai}.

\begin{table}\tcp{\tsc{Taliabu} and nearby languages *-əj}{08-06}
	\begin{threeparttable}\st{6pt}
	\begin{tabular}{lllll}\lsptoprule
*gloss	&	small worm	&	navel	&	fly (insect)	&	vine, rope\\
PMP	&	*qul\tbr{əj}	&	*pus\tbr{əj}	&	*lal\tbr{əj}	&	*waR\tbr{əj}\\	\midrule
Saluan	&	{\hr}ul\tbr{o}, ul\tbr{oo}	&	{\hr}pus\tbr{oo}	&	{\hr}lal\tbr{o}, lal\tbr{oo}	&	\cg{(putal, ueʔ)}\\
Balantak	&	{\hr}ul\tbb{e}	&	{\hr}pus\tbb{e}	&	{\hr}laal\tbb{e}	&	\cg{(luiʔ, lukut, oe)}\\
Banggai	&	{\hr}ul\tbr{oy}	&	{\hr}pus\tbr{oy}	&	\cg{(poos)}	&	{\hr}al\tbr{os}\su{†}\\
\noalign{\vskip-1.5ex}
\multicolumn{5}{@{}c@{}}{\dashedline{110mm}} \\
\noalign{\vskip-0.5ex}
Soboyo	&		&	{\hr}pus\tbb{e}	&	\cg{(moŋoŋ)}	&	{\hr}wah\tbr{o}\su{Δ}\\
Taliabu	&	{\hr}ul\tbb{e} tunu\su{‡}	&	{\hr}pus\tbb{e}	&	{\hr}lal\tbb{e} {\tl} lal\tbb{i}	&	{\hr}tuka wah\tbr{o}-c\su{ǁ}\\
Kadai	&		&		&		&	{\hr}wah\tbr{o}\su{Δ}\\	\midrule
Mangole	&		&		&	\cg{(fiŋa muya)}	&	\cg{(meu)}\\
Sanana	&	{\hr}ul\tbb{i}	&	{\hr}puh\tbb{i}	&	\cg{(fina muya)}	&	{\hr}waʔ\tbb{i}\\
Lisela	&		&		&	\cg{(ɸeŋa)}	&	{\hr}wah\tbb{e}-t\\
Buru	&	{\hr}ul\tbb{e}-t	&	{\hr}pus\tbb{e}-n	&	\cg{(feŋa)}	&	{\hr}wah\tbb{e}-t, wah\tbb{e}-n\\
Ambelau	&	{\hr}ul\tbb{e}-a	&	{\hr}fus\tbb{e}	&	\cg{(bena)}	&	{\hr}βah\tbb{e}-re\\	\midrule
Kayeli 	&		&	{\hr}hus\tbb{e}-n	&	\cg{(bena)}	&	{\hr}wal\tbb{e}-t\\
		\lspbottomrule
	\end{tabular}
		\begin{tablenotes}\tnsize
			\item[†]	{Banggai \it{alos} ‘tough creeping roots or tendrils,
								rope’ has irregular final *j > \it{s}.}
			\item[‡]	{Taliabu \it{ule tunu} ‘toxic caterpillar’.}
			\item[Δ]	{Soboyo
								and Kadai \it{waho} means ‘rattan’, which grows as a vine.}
			\item[{ǁ}]	{Taliabu \it{tuka waho-c} means ‘intestines (lit: stomach rope)’.
										[c] = voiceless palatal plosive, or affricate [tsʸ].}
		\end{tablenotes}
	\end{threeparttable}
\end{table}

\clearpage
The reflexes of PMP *-uy are not useful for subgrouping purposes,
but are nevertheless interesting,
with Mangole showing a split *-uy > \it{u,
i}, and Sanana showing a loss *-uy > {\0},
as in \trf{08-07}.

\begin{table}
	\centering\tcp{\tsc{Taliabu} and nearby languages *-uy}{08-07}
	\st{6pt}\begin{tabular}{llll}\lsptoprule
*gloss	&	fire	&	pig	&	swim\\
PMP	&	*hap\tbr{uy}	&	*bab\tbr{uy}	&	*naŋ\tbr{uy}\\ \midrule
Saluan	&	{\hr}ap\tbr{u}, ap\tbr{uu}	&	{\hr}ba\tbr{u}, ba\tbr{uu}	&	{\hr}laŋ\tbr{u}, laŋ\tbr{uu}\\
Balantak	&	{\hr}ap\tbr{u}	&	{\hr}bauʔ	&	{\hr}laŋ\tbr{u}\\
Banggai	&	\cg{(bilat)}	&	{\hr}bab\tbr{ui}	&	\cg{(kayok)}\\
\noalign{\vskip-1.5ex}
\multicolumn{4}{@{}c@{}}{\dashedline{88mm}} \\
\noalign{\vskip-0.5ex}
Soboyo	&	\cg{(utoin)}	&	\cg{(mboa)}	&	{\hr}naŋ\tbr{u}\\
Taliabu	&	\cg{(utoŋ)}	&		&	{\hr}naŋ\tbr{u}\\
Kadai	&	\cg{(uto)}	&	\cg{(mboa)}	&	{\hr}naŋ\tbr{u}\\ \midrule
Mangole	&	{\hr}ap\tbr{u}	&	{\hr}faf\tbr{i}	&	{\hr}naŋ\tbr{u}\\
Sanana	&	{\hr}ap	&	{\hr}faf	&	{\hr}nan\\
Lisela	&	\cg{(bana)}	&	{\hr}ɸaɸ\tbr{u}	&	{\hr}naŋ\tbr{o} \\
Buru	&	\cg{(bana)}	&	{\hr}faf\tbr{u}	&	{\hr}naŋ\tbr{o} \it{‘wade’}\\
Ambelau	&	{\hr}af\tbr{u}	&	{\hr}baβ\tbr{u}	&	{\hr}en-nal\tbr{o}, nan\tbr{o}\\ \midrule
Kayeli	&	{\hr}ah\tbr{u}-t	&	{\hr}bab\tbr{u}	&	{\hr}en-nan\tbr{o}\\
		\lspbottomrule
	\end{tabular}
\end{table}
